\documentclass[11pt]{article}
\usepackage[margin=1in]{geometry}
\usepackage[T1]{fontenc}
\usepackage{lmodern,microtype,mathrsfs,amsmath,amssymb,amsthm,mathtools,booktabs,array,longtable}
\usepackage{xcolor}
\definecolor{linkink}{RGB}{28,63,88}
\usepackage[colorlinks=true,linkcolor=linkink,citecolor=linkink,urlcolor=linkink]{hyperref}
\hypersetup{pdftitle={An entropy--variance profile for Boolean noise: I16.1},pdfauthor={Dimiter Kramer}}
\usepackage{fancyhdr}
\pagestyle{fancy}\fancyhf{}
\fancyhead[L]{\small Entropy profiles for Boolean noise}
\fancyhead[R]{\small I16.1: compressed certificates}
\fancyfoot[C]{\small\thepage}
\setlength{\headheight}{14pt}
\setlength{\parskip}{3pt}
\setlength{\emergencystretch}{1em}
\newtheorem{theorem}{Theorem}[section]
\newtheorem{lemma}[theorem]{Lemma}
\newtheorem{proposition}[theorem]{Proposition}
\newtheorem{corollary}[theorem]{Corollary}
\theoremstyle{remark}\newtheorem{remark}[theorem]{Remark}
\newcommand{\E}{\mathbb E}\newcommand{\PP}{\mathbb P}\newcommand{\dd}{\,\mathrm d}
\newcommand{\eps}{\varepsilon}\newcommand{\atanh}{\operatorname{atanh}}
\newcommand{\KL}{\operatorname{D}}\newcommand{\MI}{\operatorname I}
\newcommand{\Bprof}{\mathsf B_p}\newcommand{\Pprof}{\Phi_p}
\newcommand{\privA}{\mathfrak a}\newcommand{\Afun}{\mathscr A}
\numberwithin{equation}{section}
\title{An entropy--variance profile for Boolean noise}
\author{Dimiter Kramer\thanks{Developed with substantial AI assistance. This is a computer-assisted proof candidate for independent mathematical assessment, not an independently established resolution or a proof-kernel-verified theorem. The rational certificates are replayed in the supplied package; the analytic and probabilistic arguments remain subject to independent scrutiny.}}
\date{20 September 2026 \quad Unified closure, compressed certificates I16.1}
\begin{document}
\maketitle
\begin{abstract}
We give a computer-assisted candidate proof of an entropy--variance profile
for deterministic Boolean functions, uniform in finite dimension, output mean,
and binary-symmetric noise. The same positive blend of two entropy lower bounds
closes the induction at every nontrivial noise level. Its finite-gap deficit
reduces to three positive geometric coefficients and one constrained
entropy-product inequality. A differential channel estimate both controls the
compensation and makes the product target monotone. This permits a compact
certificate of 184 positive quartics and 36 entropy-infeasible boxes, with exact
rational replay; the two noise endpoints are treated analytically. The geometric
checks use 71 positive coefficients and two-term logarithm bounds.
The profile, all-noise blend, quantitative remainder, and equality cases of I16
are unchanged. This revision reduces the verification burden, not the domain.
The supplied replay checks finite certificates; independent mathematical review
and end-to-end proof-kernel verification remain outstanding.
\end{abstract}
% === 01_profile.tex ===
\section{The profile and its consequences}\label{sec:intro}
Let $X=(X_1,\ldots,X_n)$ be uniform on $\{-1,1\}^n$. Independently flip
each sign with probability $p\in[0,1/2]$ to obtain $Y$. For a deterministic
$f:\{-1,1\}^n\to\{0,1\}$, put $F=2f-1$ and $m=\E F$. All entropies
are in nats. Throughout,
\[
 h(a)=-a\ln a-(1-a)\ln(1-a),\qquad
 \eta(u)=h((1+u)/2),\qquad L=\ln2,
\]
with $0\ln0=0$. A coordinate indicator has mutual information $L-h(p)$.
The Courtade--Kumar problem asks whether any Boolean function can do better~\cite{CK}.

Write
\[
 \rho=1-2p,\quad R=\rho^2,\quad q=1-R,\quad V(m)=1-m^2,
 \qquad \mathscr E_m=\eta(m)-LV(m).
\]
The proposed lower bound separates a variance term from an entropy excess:
\begin{equation}\label{eq:profile}
 \boxed{\Bprof(m)=h(p)V(m)+q\mathscr E_m
       =q\eta(m)+\theta V(m),\qquad \theta=h(p)-qL.}
\end{equation}
\begin{samepage}\begin{theorem}[Entropy--variance profile]\label{thm:main}
For every integer $n\ge0$, every deterministic Boolean $f$, and every
$p\in[0,1/2]$,
\begin{equation}\label{eq:main-profile}
 H(f(X)\mid Y)\ge\Bprof\bigl(\E[2f(X)-1]\bigr).
\end{equation}
\end{theorem}\end{samepage}
The first form of the profile describes the bound; the second is convenient
for restriction. The coefficient $q$ depends only on the channel.

\subsection{Why the formula has the right endpoints}
Set $\kappa(u)=L-\eta(u)$. Integration of
$\eta'(u)=-\atanh u$ gives
\begin{equation}\label{eq:kappa-series}
 \kappa(u)=\sum_{j\ge1}\frac{u^{2j}}{(2j)(2j-1)},\qquad
 \frac{u^2}{2}\le\kappa(u)\le Lu^2\quad (|u|\le1).
\end{equation}
The endpoint series converges, and its coefficients sum to $L$.
Consequently $\mathscr E_m\ge0$,
$0\le\theta\le(L-1/2)R$, and
\[
 \Bprof(0)=h(p),\qquad \Bprof(\pm1)=0,\qquad
 \mathsf B_0(m)=0,\qquad \mathsf B_{1/2}(m)=\eta(m).
\]
These correspond to coordinates, constants, zero noise, and complete noise.
Also $\Bprof$ is even and concave, and
\[
 \eta(m)-\Bprof(m)=R\mathscr E_m+\kappa(\rho)V(m)\ge0.
\]
Thus the local entropy budgets used in the proof are feasible.

\begin{corollary}[Courtade--Kumar consequence]\label{cor:CK}
Theorem~\ref{thm:main} implies, on the same full domain,
\begin{equation}\label{eq:CK}
 \MI(f(X);Y)\le L-h(p)=\kappa(\rho).
\end{equation}
\end{corollary}
\begin{proof}
Since $H(f(X))=\eta(m)$, it suffices to observe that
\begin{equation}\label{eq:bias-payment}
 \Bprof(m)-\eta(m)+\kappa(\rho)
 =R\kappa(m)-\theta m^2\ge(1-L)Rm^2\ge0.
\end{equation}
Here $\kappa(m)\ge m^2/2$, $\theta\le(L-1/2)R$, and $L<1$.
\end{proof}
For $0<p<1/2$, Section~\ref{sec:equality} derives equality only for
constants and signed coordinates in the profile, and only for signed
coordinates in~\eqref{eq:CK}. A signed coordinate means
$f(x)=\mathbf1\{\sigma x_i=1\}$, $\sigma\in\{-1,1\}$.


\subsection{One closure mechanism, and its verification boundary}
The proof spine is unchanged: two all-noise lower bounds, one compensated
blend, and a positive finite-gap remainder. The auxiliary product lemma is
proved by analytic caps and a rational compact certificate. The latter now
fits in six lines, reproduced in Appendix~\ref{app:replay}; every leaf is
verified on its whole box, not sampled. Numerical partitions remain inside
this subsidiary verification, not in the choice of closure mechanism.

\begin{center}
\begin{tabular}{p{.25\textwidth}p{.57\textwidth}p{.09\textwidth}}
\toprule
Ingredient & Role & Section\\
\midrule
Quarter-variance induction & Two-bit base and a coordinate supplied in every higher dimension & 2\\
Two all-noise bounds & Prior-sensitive consistency and sharp entropy/gap bound & 3--4\\
One compensated blend & Exact finite-gap cancellation; no extremizer hypothesis & 5\\
Scalar ingredients & Channel compensation, certified product, geometric coefficients & 6--8\\
Equality & Strict closure, or strict entropy concavity at zero gap & 9\\
\bottomrule
\end{tabular}
\end{center}
% === 02_induction.tex ===
\section{A quarter-variance gap and a two-bit base}\label{sec:induction}
For posterior probabilities define
\[
 J(a,b)=\frac{h(a)+h(b)}2,\qquad
 M_p(a,b)=\frac{h((1-p)a+pb)+h(pa+(1-p)b)}2.
\]
The local theorem keeps the actual means and is used only on the domain
needed by the induction.
\begin{theorem}[Posterior closure at a quarter-variance gap]\label{thm:closure}
Fix $p\in[0,1/2]$. Let $A,B\in[0,1]$ be random variables on one finite
probability space, and put
\[
 \mu_0=\E A,\quad\mu_1=\E B,\quad
 m=\mu_0+\mu_1-1,\quad d=|\mu_1-\mu_0|.
\]
Suppose $4d^2\le1-m^2$ and
\begin{equation}\label{eq:childhyp}
 \E J(A,B)\ge\frac{\Bprof(2\mu_0-1)+\Bprof(2\mu_1-1)}2.
\end{equation}
Then
\begin{equation}\label{eq:closure}
 \E M_p(A,B)\ge\Bprof(m).
\end{equation}
No independence of $A,B$ is assumed.
\end{theorem}
Section~\ref{sec:blend} proves this statement with one all-noise blend. The next lemma
explains why the gap hypothesis does not restrict the final theorem.

\begin{lemma}[A quarter-variance coordinate]\label{lem:normalized-gap}
Let $F:\{-1,1\}^n\to\{-1,1\}$, with $n\ge3$, and let $m=\E F$.
Among any three distinct coordinates, at least one satisfies
\begin{equation}\label{eq:normalized-gap}
 d:=|\E(FX_i)|\le\frac12,\qquad 4d^2\le1-m^2.
\end{equation}
\end{lemma}
\begin{proof}
Complement $F$ so that $m\ge0$, and choose signs so that
$T=\sigma_iX_i+\sigma_jX_j+\sigma_kX_k$ satisfies
$\E(FT)=|\E(FX_i)|+|\E(FX_j)|+|\E(FX_k)|$.
Independence and fairness give $\E|T+1|=3/2$ and $\E T=0$.
The bounds $F(T+1)\le|T+1|$ and $F(T+3)\le T+3$ yield
\[
 \E(FT)\le\min\{3/2-m,\,3(1-m)\}.
\]
For $0\le m\le3/4$ and $3/4\le m\le1$, respectively,
\begin{align*}
 \frac94(1-m^2)-(3/2-m)^2
 &=\frac{m(12-13m)}4\ge0,\\
 \frac94(1-m^2)-9(1-m)^2
 &=\frac94(1-m)(5m-3)\ge0.
\end{align*}
Thus the smallest correlation obeys
$3d\le\E(FT)\le\tfrac32\sqrt{1-m^2}$.
This gives $4d^2\le1-m^2$ and hence $d\le1/2$.
\end{proof}
The bound is sharp: three-bit majority has mean zero and all three
coordinate correlations equal to $1/2$.
A two-bit conjunction does not satisfy this condition at either coordinate,
which is why the base is proved separately.

\begin{lemma}[The two-bit base]\label{lem:two-bit}
The profile inequality holds for every Boolean function of at most two
variables. For $0<p<1/2$, the only profile equality cases in this class
are constants and signed coordinates.
\end{lemma}
\begin{proof}
Under input permutations, input sign changes, and output complementation,
every two-bit Boolean function is a constant, a coordinate, parity, or
conjunction. A truth table with one or three ones is in the conjunction
orbit; with two ones, the marked vertices are adjacent (a coordinate) or
opposite (parity). These operations preserve the conditional entropy and the even
profile. Constants and coordinates attain the profile. Parity has mean zero
and conditional entropy $\eta(\rho^2)>\eta(\rho)=h(p)$ at nontrivial noise.

For conjunction define its profile deficit
\begin{equation}\label{eq:two-bit-deficit}
 G(p)=\frac{h(p^2)+2h(p(1-p))+h((1-p)^2)}4
       -\frac34h(p)-4p(1-p)
        \left[h(1/4)-\frac34L\right].
\end{equation}
The four posterior probabilities give the first term directly. Thus
$G(0)=G(1/2)=0$; symmetry under $p\mapsto1-p$ gives $G'(1/2)=0$.
Writing $v=p(1-p)$, differentiation and rational simplification give
\begin{equation}\label{eq:two-bit-third}
 G'''(p)=
 \frac{(1-2p)(5v^4+2v^3+39v^2+4v+4)}
 {4v^2(2+v)^2(1-v)^2}>0\qquad(0<p<1/2).
\end{equation}
To see that these signs suffice, fix $x\in(0,1/2)$ and put
$c=G(x)/[x(x-1/2)^2]$.
The function $G(t)-ct(t-1/2)^2$ vanishes at $0,x,1/2$, and its derivative
vanishes at $1/2$. Three applications of Rolle's theorem, all derivative
points lying in $(0,1/2)$, give $6c=G'''(\xi)>0$ for some $\xi$.
Hence $G(x)>0$. Only continuity, not differentiability, at $p=0$ is needed.
The noise endpoints are immediate.
\end{proof}

\begin{proof}[Proof of Theorem~\ref{thm:main}, using Theorem~\ref{thm:closure}]
Fix $p$ and induct on $n$. Lemma~\ref{lem:two-bit} supplies $n\le2$.
For $n\ge3$, choose $i$ by Lemma~\ref{lem:normalized-gap}, write $f_0,f_1$
for the two sections, and let $W$ be the noisy observation of the other
coordinates. Define $A=\E[f_0\mid W]$ and $B=\E[f_1\mid W]$.
Then $m=\E[2f-1]$ and
$d=|\E f_1-\E f_0|=|\E((2f-1)X_i)|$, so the local quarter-variance condition holds.
The two induction hypotheses imply~\eqref{eq:childhyp}.
The selected noisy sign is fair and independent of $W$, with posterior
input probabilities $1-p,p$. Its two output posteriors are
$(1-p)A+pB$ and $pA+(1-p)B$, each with conditional probability $1/2$.
Therefore $H(f\mid Y)=\E M_p(A,B)$, and~\eqref{eq:closure} proves the parent bound.
The endpoints $p=0,1/2$ follow directly as well. Corollary~\ref{cor:CK}
gives the information bound.
\end{proof}

% === 03_recombination.tex ===
\section{Mean consistency and entropy recombination}\label{sec:privacy}
Let $S$ be a uniform sign independent of a finite random variable $W$, and let $Z$ be a binary random variable. It may be randomized given $(S,W)$. Write $S_\rho=SN_S$, where the noise sign $N_S$ is independent of $(S,W,Z)$ and
$\PP(N_S=-1)=p$.
Although $S$ and $W$ are independent, they need not remain independent after $Z$ is known. We retain that conditional dependence instead of discarding it.

\subsection{Two elementary relative-entropy estimates}
For Bernoulli probabilities let
\[
\KL(x\Vert y)=x\ln(x/y)+(1-x)\ln((1-x)/(1-y)).
\]
For sign means, the same divergence is the Bregman divergence of $\kappa$; its second derivative is $\kappa''(u)=(1-u^2)^{-1}$. Consequently
\begin{equation}\label{eq:binary-pinsker}
\KL\bigl((1+u)/2\Vert(1+v)/2\bigr)\ge\tfrac12(u-v)^2.
\end{equation}
Also the second derivative of $u\mapsto\kappa(\rho u)$ is
$R/(1-Ru^2)\le R/(1-u^2)$. Integrating the second derivatives along a line segment gives
\begin{equation}\label{eq:onebit-contraction}
\KL\bigl((1+\rho u)/2\Vert(1+\rho v)/2\bigr)
\le R\KL\bigl((1+u)/2\Vert(1+v)/2\bigr).
\end{equation}
This is a \emph{one-bit} statement for arbitrary priors. We do not assert the corresponding product-channel mutual-information contraction for arbitrary input laws. Averaging \eqref{eq:onebit-contraction} over conditional laws gives
\begin{equation}\label{eq:conditional-contraction}
\MI(S_\rho;W\mid Z)\le R\MI(S;W\mid Z).
\end{equation}
Endpoint cases follow by limits, or directly on deterministic supports.

We also need an upper bound with a fixed Bernoulli prior. For $1/2\le\mu<1$,
\begin{equation}\label{eq:binary-upper}
\KL(c\Vert\mu)\le\frac{-\ln\mu}{(1-\mu)^2}(c-\mu)^2,
\qquad 0\le c\le1.
\end{equation}
To prove the upper bound, observe that the continuous extension of
$\KL(c\Vert\mu)/(c-\mu)^2$ is
\[
\int_0^1\frac{1-t}{[\mu+t(c-\mu)][1-\mu-t(c-\mu)]}\dd t.
\]
It is convex as a function of $c$, since $1/[x(1-x)]=1/x+1/(1-x)$ is convex. Its maximum is therefore at $c=0$ or $c=1$. The second endpoint is larger when $\mu\ge1/2$: the difference
$-\mu^2\ln\mu+(1-\mu)^2\ln(1-\mu)$ has derivative $2h(\mu)-1$, is concave on $[1/2,1]$, and vanishes at both ends. This proves \eqref{eq:binary-upper}.

\subsection{Dependence forced by the actual mean gap}
Orient the binary variables so that
\[
 \mu=\PP(Z=1)\in[1/2,1),\qquad
 d=\PP(Z=1\mid S=1)-\PP(Z=1\mid S=-1)\ge0,
 \qquad m=2\mu-1.
\]
Write $V=1-m^2$ and
\begin{equation}\label{eq:a}
 a(m)=2(1+m)^2\ln\frac2{1+m}=8\mu^2(-\ln\mu).
\end{equation}
The bounds $1-\mu\le-\ln\mu\le(1-\mu)/\mu$ imply
\begin{equation}\label{eq:a-bounds}
 V\le a(m)\le2V.
\end{equation}
Set $I_W=I(Z;W)$, $T_W=I(S;W\mid Z)$, and
$D_W=I(Z;W\mid S)=I_W+T_W$, the last equality using $S\perp W$.

\begin{lemma}[Conditional dependence forced by a mean gap]\label{lem:dependence}
In this finite experiment,
\begin{equation}\label{eq:privacy-ratio}
 T_W\ge\frac{d^2}{a(m)}I_W,
 \qquad
 T_W\ge\frac{d^2}{a(m)+d^2}D_W.
\end{equation}
\end{lemma}
\begin{proof}
Only $d>0$ requires proof. Let $c=\PP(Z=1\mid W)$ and let $t$ be the
corresponding difference between the two section probabilities. Then
$\E c=\mu$, $\E t=d$. Put $b=\mu(1-\mu)$ and $v=\operatorname{Var}(c)$.
Conditioning on $Z$ gives posterior sign means $t/(2c)$ and
$-t/[2(1-c)]$, with prior means $d/(2\mu)$ and $-d/[2(1-\mu)]$.
Apply \eqref{eq:binary-pinsker} and expand the resulting two variances:
\[
 T_W\ge\frac18\left(\E\frac{t^2}{c(1-c)}-\frac{d^2}{b}\right)
 \ge\frac{d^2v}{8b(b-v)}\ge\frac{d^2v}{8b^2}.
\]
The middle inequality is weighted Cauchy--Schwarz, since
$\E[c(1-c)]=b-v$. This denominator is positive: if it vanished, $c$ would
be zero or one almost surely, forcing $t=0$ and contradicting $d>0$.
At such individual endpoints the ratio is defined as zero because $t=0$.
On the other hand, \eqref{eq:binary-upper} gives
\[
 I_W=\E D(c\Vert\mu)\le\frac{-\ln\mu}{(1-\mu)^2}v.
\]
Multiplication by $d^2/a(m)$ proves the first assertion without division by
$v$. The identity $D_W=I_W+T_W$ gives the second. For $d=0$, both follow
from nonnegativity.
\end{proof}

\subsection{The recombination identity}
Define
\[
\mathcal H(m,d)=\frac{\eta(m+d)+\eta(m-d)}2,\qquad
\bar J=\E J(A,B),\quad \bar M=\E M_p(A,B),
\]
where $A=\PP(Z=1\mid S=-1,W)$ and $B=\PP(Z=1\mid S=1,W)$. Set
\[
J_*=\mathcal H(m,d)=H(Z\mid S),\qquad M_*=\mathcal H(m,\rho d)=H(Z\mid S_\rho),\qquad
\gamma=\frac{q d^2}{\privA(m)+d^2}.
\]
\begin{theorem}[Mean-sensitive recombination]\label{thm:privacy}
For $0\le m<1$, $0\le d\le1-m$, and $0<p<1/2$,
\begin{equation}\label{eq:privacy-bound}
\bar M\ge M_*-(1-\gamma)(J_*-\bar J).
\end{equation}
The theorem holds for every joint posterior law, with the means stated above. Complementation and exchange of the two sections remove the sign conventions on $m,d$.
\end{theorem}
\begin{proof}
The chain rule and $S\perp W$, $S_\rho\perp W$ give
\begin{align*}
(\bar M-\bar J)-(M_*-J_*)
&=\MI(Z;W\mid S)-\MI(Z;W\mid S_\rho)\\
&=\MI(S;W\mid Z)-\MI(S_\rho;W\mid Z)\\
&\ge q\MI(S;W\mid Z)\\
&\ge\gamma\MI(Z;W\mid S)=\gamma(J_*-\bar J).
\end{align*}
The last two inequalities are \eqref{eq:conditional-contraction} and \eqref{eq:privacy-ratio}. Rearrangement gives \eqref{eq:privacy-bound}.
\end{proof}
This bound retains the prior midpoint through $\privA(m),J_*,M_*$. That information is absent from an estimate expressed only in terms of average child entropy and an absolute posterior difference.



Every finite joint law of the posterior pair $(A,B)$ is covered by this
experiment: take $W$ to be its underlying finite outcome, $S$ independently
fair, and $\PP(Z=1\mid S=-1,W)=A(W)$, $\PP(Z=1\mid S=1,W)=B(W)$.
For an entropy lower budget $\bar J\ge j$, the coefficient $1-\gamma$ is
nonnegative, so Theorem~\ref{thm:privacy} also gives
$\bar M\ge M_*-(1-\gamma)(J_*-j)$.
\section{The sharp entropy--gap bound}\label{sec:budget}
This bound retains the average child entropy before estimating the mixing
gain. Its sharpness fixes the entropy budget and the mean gap, not an
additional parent mean. No mean-preserving extremizer claim is used below.

\subsection{Centering at fixed entropy per unit gap}
For $0\le\rho\le1$ and $e\ge0$, let $\tau(e)\in(0,1]$ be the unique solution of
\[
 \frac{\eta(\tau(e))}{\tau(e)}=e,
 \qquad \mathcal F_\rho(e)=\frac{\eta(\rho\tau(e))}{\tau(e)}.
\]
Existence and uniqueness follow from
\begin{equation}\label{eq:I15-entropy-ratio}
 \left(\frac{\eta(t)}t\right)'=-\frac{\mathscr A(t)}{t^2}<0,
 \qquad \mathscr A(t)=\eta(t)-t\eta'(t)
       =L-\frac12\ln(1-t^2),
\end{equation}
and the limits $\eta(t)/t\to\infty$ at zero and $\eta(1)=0$.

\begin{lemma}[Fixed-budget centering]\label{lem:I15-perspective}
For $a,b\in[0,1]$ with $t=|a-b|>0$,
\begin{equation}\label{eq:I15-pointwise-budget}
 \frac{M_p(a,b)}t\ge
 \mathcal F_\rho\!\left(\frac{J(a,b)}t\right),\qquad 0\le p\le1/2.
\end{equation}
\end{lemma}
\begin{proof}
The cases $p=0$, $a+b=1$, and $J(a,b)=0$ are immediate. Otherwise write
$x=|a+b-1|>0$ and $\mathcal H(x,t)=[\eta(x+t)+\eta(x-t)]/2$.
In the interior $x+t<1$, put
\[
 \ell=-\mathcal H_x>0,\quad \jmath=-\mathcal H_t\ge0,\quad
 Q=\mathcal H+t\jmath>0,\quad P=Q+x\ell,\quad
 R_0=\frac{1+x^2-t^2}{2x},\quad \Theta=P-R_0\ell.
\]
Direct differentiation gives
\begin{equation}\label{eq:I15-ratio-calculus}
 \ell_t=\frac{(1-(x+t)^2)^{-1}-(1-(x-t)^2)^{-1}}2>0,
 \quad P_t=R_0\ell_t,\quad \Theta_t=\frac tx\ell.
\end{equation}
At $t=0$,
\[
 \Theta(x,0)=\frac{\Omega(x)}{2x},\qquad
 \Omega(u)=2u\eta(u)-(1-u^2)\atanh u.
\]
The endpoint limits are $\Omega(0)=\Omega(1)=0$ and
$\Omega''(u)=-2\atanh u\le0$. Thus $\Theta\ge0$ and
\begin{equation}\label{eq:I15-decreasing-ratio}
 \partial_t(Q/\ell)=\partial_t(P/\ell)
   =-\frac{\Theta\ell_t}{\ell^2}\le0.
\end{equation}

Fix the actual ratio $e=\mathcal H(x,t)/t>0$ and follow the level curve
$\mathcal H(x,t(x))/t(x)=e$ toward $x=0$.
This curve exists on the whole intervening interval: at each fixed $x$ the
ratio decreases strictly in $t$, from infinity to
$h(x)/[2(1-x)]$ at $t=1-x$, and this last value increases with $x$, since
its derivative is $-\ln x/[2(1-x)^2]>0$.
Along the interior of the curve,
\begin{equation}\label{eq:I15-level-derivative}
 t'(x)=-\frac{t\ell(x,t)}{Q(x,t)},\qquad
 \frac{\mathrm d}{\mathrm dx}\frac{\mathcal H(x,\rho t(x))}{t(x)}
 =\frac1t\left[
 \frac{\ell(x,t)Q(x,\rho t)}{Q(x,t)}-\ell(x,\rho t)\right]\ge0.
\end{equation}
The sign follows from~\eqref{eq:I15-decreasing-ratio} because $\rho t\le t$.
At $x=0$ the curve has $t(0)=\tau(e)$, proving the claim.
Continuity supplies $x+t=1$; when $J=0$ and $t>0$, the pair is $(0,1)$
up to exchange and the formula is an equality. No probability law has
been changed in this pointwise comparison.
\end{proof}

\subsection{The sharp bound determined by the child budget}
\begin{theorem}[Entropy-budget bound]\label{thm:I15-budget}
Let $A,B\in[0,1]$ be random variables on one finite probability space,
let $d=|\E A-\E B|>0$, and suppose
$\E J(A,B)\ge j$ with $0\le j\le\eta(d)$. Define $\tau\in[d,1]$ by
\begin{equation}\label{eq:I15-budget-root}
 \frac{\eta(\tau)}\tau=\frac jd.
\end{equation}
Then, for every $0\le p\le1/2$,
\begin{equation}\label{eq:I15-budget-bound}
 \boxed{\E M_p(A,B)\ge\mathcal C_p(j,d):=
                    \frac d\tau\eta(\rho\tau).}
\end{equation}
The bound is sharp under these two constraints. No fixed parent mean is
included in this sharpness assertion. If $d=0$, the bound used is simply
$\E M_p(A,B)\ge j$.
\end{theorem}
\begin{proof}
For $p=0$ the assertion is $\E J\ge j$, so assume $0\le\rho<1$.
The function $\mathcal F_\rho$ is increasing and convex. Indeed, for $0\le\rho<1$
and $e>0$, the chain rule gives
\begin{equation}\label{eq:I15-budget-slope}
 \mathcal F_\rho'(e)=\frac{\mathscr A(\rho\tau(e))}{\mathscr A(\tau(e))}
                  \in(0,1).
\end{equation}
The ratio decreases with $\tau$: write
$\mathscr A(\sqrt x)=\sum_{k\ge0}b_kx^k$ with
$b_0=L$ and $b_k=1/(2k)>0$ for $k\ge1$.
After differentiation, the numerator of the derivative of
$\mathscr A(\rho\sqrt x)/\mathscr A(\sqrt x)$ is
\[
 \sum_{i>k\ge0}(i-k)b_i b_k(\rho^{2i}-\rho^{2k})x^{i+k-1}\le0.
\]
Here $\rho^0=1$, including when $\rho=0$.
Absolute convergence justifies the calculation for $0<x<1$.
Since $\tau(e)$ decreases, the slope in~\eqref{eq:I15-budget-slope}
increases with $e>0$. Finally $\tau(e)\to1$ as $e\downarrow0$,
so $\mathcal F_\rho(e)\to h(p)$; continuity extends convexity to $e=0$.
No derivative of $\mathscr A$ at $t=1$ is needed.

For $j>0$ put $e_*=j/d$, $\lambda=\mathcal F_\rho'(e_*)$ and
$c=\mathcal F_\rho(e_*)-\lambda e_*$.
One has $c>0$: in its expression
$[\eta(\rho\tau)-\lambda\eta(\tau)]/\tau$, both
$\eta(\rho\tau)\ge\eta(\tau)$ and $\lambda<1$ apply.
A supporting line to the convex function, followed by Lemma~\ref{lem:I15-perspective},
gives pointwise
\[
 M_p(a,b)\ge\lambda J(a,b)+c|a-b|.
\]
At $a=b$ this follows instead from $M_p=J$ and $\lambda\le1$.
Averaging and using $\E|A-B|\ge d$ proves
$\E M_p\ge\lambda j+cd=\mathcal C_p(j,d)$.
For $j=0$, monotonicity gives $\mathcal F_\rho(e)\ge\mathcal F_\rho(0)=h(p)$
and the same argument applies with $\lambda=0,c=h(p)$.

For sharpness, put $w=d/\tau\le1$. Assign mass $w$ to the centered pair
$((1-\tau)/2,(1+\tau)/2)$ and mass $(1-w)/2$ to each of $(0,0)$ and $(1,1)$.
Then the mean difference is $d$, the average child entropy is
$w\eta(\tau)=j$, and the mixed entropy is $w\eta(\rho\tau)$.
The case $d=0$ follows from entropy concavity.
\end{proof}


\section{One compensated blend on the entire noise range}\label{sec:blend}
Throughout Sections 5--8, use the nondegenerate domain
\begin{equation}\label{eq:I16-domain}
 0<R<1,\qquad 0\le m<1,\qquad
 0<d\le\min\{1-m,\sqrt{1-m^2}/2\}.
\end{equation}
Set $V=1-m^2$, $s=d^2$, $H=\eta(\rho)=h(p)$ and
\[
 q=1-R,\qquad \theta=H-qL,\qquad
 \alpha=\theta/q,\qquad u=\alpha/R,\qquad
 \Omega=\eta(m)+\alpha V,\qquad B=\Bprof(m)=q\Omega.
\]
The symbol $\alpha$ here is a channel parameter; the mixing coefficient of
Section 3 was denoted $\gamma$. Write
\begin{equation}\label{eq:I16-budget}
 J_*=\mathcal H(m,d),\quad M_*=\mathcal H(m,\rho d),\quad
 j=qJ_*+\theta(V-s),\quad a=\privA(m),\quad A=a/V,\quad U=1+2\alpha V.
\end{equation}
For the prescribed means and $\bar J\ge j$, Sections 3--4 give
\begin{equation}\label{eq:I16-two-bounds}
 \bar M\ge C:=\frac d\tau\eta(\rho\tau),\quad
 \frac{\eta(\tau)}\tau=\frac jd;\qquad
 \bar M\ge P:=M_*-\left(1-\frac{qs}{a+s}\right)(J_*-j).
\end{equation}
Here $0<j\le J_*\le\eta(d)$ and $d\le\tau<1$. The lower strictness of
$j$ follows from $j\ge H(V-s)>0$; the upper bounds follow from
$\Bprof\le\eta$ and concavity of $\eta$.

Choose the positive weights
\begin{equation}\label{eq:I16-weights}
 \boxed{w_C=\frac{2(C+B-j)}q,\qquad
 w_P=\frac aV\frac{1+2\alpha V}{R}.}
\end{equation}
They are positive because $C\ge j$, $B>0$, and the remaining factors are
positive. In particular, the denominator $a+s$ inside $P$ is retained.

\begin{theorem}[All-noise compensated closure]\label{thm:I16-blend}
On \eqref{eq:I16-domain}, define
$\mathcal W=w_C(C-B)+w_P(P-B)$. Then
\begin{equation}\label{eq:I16-margin}
 \boxed{\mathcal W>
 qd^2\left(\frac1{250}+\frac\alpha{10}+\frac{\alpha^2}{53}\right)>0.}
\end{equation}
Consequently the same convex blend gives, for every admissible posterior law,
\begin{equation}\label{eq:I16-assembly}
 \bar M\ge\frac{w_CC+w_PP}{w_C+w_P}
 =B+\frac{\mathcal W}{w_C+w_P}>B.
\end{equation}
The theorem uses the certified scalar product lemma in Section 7 and the
rational geometric certificates in Section 8.
\end{theorem}

\subsection{The exact finite-gap identity}
Define
\[
 D=\eta(m)-J_*,\qquad \overline D=D/s,\qquad
 c_k=\frac{(1-m)^{1-2k}+(1+m)^{1-2k}}{4k(2k-1)}.
\]
Expanding the exact binary-information identity
\begin{equation}\label{eq:I16-D}
 D=\frac{1+m}{2}\kappa\!\left(\frac d{1+m}\right)
  +\frac{1-m}{2}\kappa\!\left(\frac d{1-m}\right)
  =\sum_{k\ge1}c_ks^k
\end{equation}
gives $c_1=1/(2V)$, $c_2=(4-3V)/(12V^3)$. The series converges at
$d=1-m$, by the convergent positive entropy series. Put
$G_n(R)=1+R+\cdots+R^n$. Direct substitution in $P$ yields
\begin{align}
 P-B&=\frac{Rq s}{a+s}\,E,\label{eq:I16-P-E}\\
 E&=\eta(m)-u(a+qV+Rs)+s(ac_2-c_1)\notag\\
 &\quad+\sum_{k\ge3}c_k
       \left[aG_{k-2}(R)s^{k-1}+RG_{k-3}(R)s^k\right].\label{eq:I16-E}
\end{align}
No child slack has been set to zero in the posterior law: $j$ is a lower
budget inserted into the two increasing lower bounds.

Let
\begin{equation}\label{eq:I16-E-L}
 E_{\rm b}=\eta(m)+a\overline D-\frac{a+s}{2V},\qquad
 E_L=E_{\rm b}-u(a+qV+Rs).
\end{equation}
The discarded series has an explicit sign:
\begin{equation}\label{eq:I16-E-remainder}
 E-E_L=R(a+s)\sum_{k\ge3}c_kG_{k-3}(R)s^{k-1}\ge0.
\end{equation}
The entropy-product kernel is
\begin{equation}\label{eq:I16-kernel}
 \mathcal K_R(t)=
 \frac{2\eta(\rho t)[\eta(\rho t)-\eta(t)]}{q t^2}.
\end{equation}
Since $j=d\eta(\tau)/\tau$ and
$B-j=q(D+\alpha s)$, the identity
$(C-B)(C+B-j)=C(C-j)-B(B-j)$ gives
\begin{equation}\label{eq:I16-exact-W}
 \boxed{\frac{\mathcal W}{s}
 =\mathcal K_R(\tau)-2q\Omega(\overline D+\alpha)
 +\frac{qaU}{V(a+s)}E.}
\end{equation}
This identity, rather than a truncated small-gap expansion, carries the
compensation through finite gaps.

\subsection{Proof of the positive remainder}
Feasibility and quarter variance give $V\ge d(2-d)$ and $V\ge4d^2$.
Their exact convex combination is
\[
 V-d^2-\frac65d
 =\frac35\{V-d(2-d)\}+\frac25\{V-4d^2\}\ge0.
\]
Thus
\begin{equation}\label{eq:I16-beta}
 \beta=\min\left\{\frac{V-s}{d},2\right\}\in[6/5,2],\qquad
 \frac{\eta(\tau)}\tau=\frac jd\ge\beta H.
\end{equation}
Lemma~\ref{lem:I16-product} applies and gives
\[
 \mathcal K_R(\tau)\ge
 q\left[L+\frac{13}{4}\alpha+2(\beta+1)\alpha^2\right].
\]
Write $c_0=L-1/2$. The channel estimate in Lemma~\ref{lem:I16-channel}
implies
\begin{equation}\label{eq:I16-comp-use}
 (qu)U=\frac\theta R(1+2\alpha V)\le c_0.
\end{equation}
Since $u=\alpha+qu$, equation~\eqref{eq:I16-E-L} can be written as
\[
 E_L=E_{\rm b}-\alpha(a+s)-(qu)(a+V).
\]
Insert these estimates into \eqref{eq:I16-exact-W} and collect powers of
$\alpha$. All multiplying factors are positive before each substitution.
The result is
\begin{equation}\label{eq:I16-quadratic}
 \frac{\mathcal W}{q s}\ge F_0+F_1\alpha+2G\alpha^2,
\end{equation}
where
\begin{align}
 F_0&=L-2\eta(m)\overline D+
 \frac a{V(a+s)}[E_{\rm b}-c_0(a+V)],\label{eq:I16-F0}\\
 F_1&=\frac{13}{4}-2\eta(m)-2V\overline D
        +\frac{2a}{a+s}E_{\rm b}-\frac aV,\label{eq:I16-F1}\\
 G&=\beta+1-a-V.\label{eq:I16-G}
\end{align}
Lemma~\ref{lem:I16-geometry} proves $F_0>1/250$, $F_1>1/10$, and
$G>1/106$ over the whole parent domain. Substitution proves
\eqref{eq:I16-margin}; averaging the two bounds with their positive weights
proves \eqref{eq:I16-assembly}. This proves Theorem~\ref{thm:I16-blend}.

\paragraph{Completion of Theorem~\ref{thm:closure}.}
Complementation and section exchange orient the means as above. At $d=0$,
the child means agree, so entropy concavity gives $\bar M\ge\bar J\ge\Bprof(m)$.
For deterministic output both posteriors are deterministic. At $p=0$ the
profile is zero. At $p=1/2$, the child budget and Jensen's inequality force
both posteriors to equal their means almost surely; the parent mixed entropy
is then $\eta(m)=\Bprof(m)$. For every other case, the strict compensated
blend applies. The local theorem therefore holds on the full stated domain,
and the induction of Section 2 applies without a numerical closure seam.

\section{One differential channel estimate}\label{sec:channel}
Put $c_0=L-1/2$ and retain $\alpha=\theta/q$, $q=1-R$.

\begin{lemma}[Channel compensation]\label{lem:I16-channel}
On $0<R<1$, $\alpha>0$ is increasing and
\begin{equation}\label{eq:I161-differential}
 \boxed{q\alpha'(R)<c_0+\frac{\alpha(R)}6.}
\end{equation}
Consequently
\begin{equation}\label{eq:I16-channel}
 \boxed{\theta\left(1+\frac{2\theta}{q}\right)
 <c_0R-\frac{R^2}{200}.}
\end{equation}
\end{lemma}
\begin{proof}
The entropy series gives
\begin{equation}\label{eq:I16-u-series}
 \alpha(R)=\sum_{j\ge0}\tau_jR^{j+1},\qquad
 \tau_j=\sum_{k\ge j+2}\frac1{(2k)(2k-1)}
       =\int_0^1\frac{x^{2j+2}}{1+x}\,\mathrm dx>0.
\end{equation}
Thus $\tau_0=c_0$, $\alpha\ge c_0R$, and $\alpha'>0$.
Write $a_{n+1}=1/[(2n+2)(2n+1)]$. The coefficient of $R^n$, $n\ge1$,
in $c_0+\alpha/6-q\alpha'$ is
\[
 \frac{(6n+1)a_{n+1}-5\tau_n}{6}.
\]
At $n=1$ its numerator is $7/2-5L>0$. For $n\ge2$, integration of
$(1+x)^{-1}\le(2x)^{-1}$ gives $\tau_n\le1/[4(n+1)]$, so the numerator
is at least
\[
 \frac{6n+1}{(2n+2)(2n+1)}-\frac5{4(n+1)}
 =\frac{2n-3}{4(n+1)(2n+1)}>0.
\]
Termwise differentiation is justified on each compact subinterval of $R<1$;
this proves \eqref{eq:I161-differential} for the entire open interval.

For $Z=q(\alpha+2\alpha^2)=\theta(1+2\theta/q)$, it follows that
\[
 Z'=(1+4\alpha)q\alpha'-\alpha-2\alpha^2
 <c_0-\left(\frac56-4c_0\right)\alpha-\frac43\alpha^2.
\]
Using $\alpha\ge c_0R$ and $Z(0)=0$, integration gives
\[
 Z<c_0R-\frac{c_0}{2}\left(\frac56-4c_0\right)R^2.
\]
The simple bounds $693/1000<L<347/500$ imply
$c_0(5/6-4c_0)/2>8299/1500000>1/200$, proving the conclusion.
\end{proof}

\begin{corollary}[The product target is monotone]\label{cor:I161-target}
For fixed $0\le\gamma\le3$, set
\[
 T_\gamma=q\left[L+\frac{13}{4}\alpha+2\gamma\alpha^2\right].
\]
This function decreases with $R$, hence increases with $p$, and satisfies
$T_\gamma<L$ at nontrivial noise.
\end{corollary}
\begin{proof}
Using \eqref{eq:I161-differential} and $L<7/10$,
\begin{align*}
 \frac{\mathrm dT_\gamma}{\mathrm dR}
 &<\frac{13}{4}c_0-L
   +\left(4\gamma c_0-\frac{65}{24}\right)\alpha
   -\frac{4\gamma}{3}\alpha^2\\
 &<-\frac1{20}-\frac{37}{120}\alpha-\frac{4\gamma}{3}\alpha^2<0.
\end{align*}
Finally $T_\gamma(0)=L$ and $R=(1-2p)^2$ decreases with $p$.
\end{proof}

\section{The certified two-variable entropy-product inequality}\label{sec:product}
\begin{lemma}[Single product bound]\label{lem:I16-product}
Let $0<p<1/2$, $0<t<1$, $6/5\le\beta\le2$, and
$\eta(t)/t\ge\beta H$. Then
\begin{equation}\label{eq:I16-product}
 \boxed{\mathcal K_R(t)\ge
 q\left[L+\frac{13}{4}\alpha+2(\beta+1)\alpha^2\right].}
\end{equation}
\end{lemma}
The coefficient, expression, and entropy constraint in this lemma are the
same throughout the proof. Analytic caps and a compact rational certificate
verify this one inequality; they do not select different posterior closure
bounds.

\subsection{A lower product with no implicit parameter}
Put $b=1-p$, $S(z)=\atanh(z)/z$ with $S(0)=1$, and
$Y=\atanh(bt)/b=tS(bt)$. Convexity of $\atanh$ gives
\[
 \eta(\rho t)-\eta(t)
 =t\int_\rho^1\atanh(vt)\,\mathrm dv
 \ge2pt\atanh(bt).
\]
Substitution in \eqref{eq:I16-kernel} gives
\begin{equation}\label{eq:I16-midpoint}
 \mathcal K_R(t)\ge K_{\rm mid}(p,t)
 :=\frac{\eta(t)}tY+\frac q2Y^2
 =\eta(t)S(bt)+\frac q2t^2S(bt)^2.
\end{equation}
The right side has a continuous value $L$ at $t=0$.
The target in Lemma~\ref{lem:I16-product} is $T_{\beta+1}$, hence
\begin{equation}\label{eq:I16-target-L}
 q\left[L+\frac{13}{4}\alpha+2(\beta+1)\alpha^2\right]<L
\end{equation}
by Corollary~\ref{cor:I161-target}.

\subsection{The complete-noise cap}
For $9/20\le p<1/2$, the entropy constraint implies $t<5/8$.
Indeed $H\ge\eta(1/10)$ and the exact logarithm bounds certify
\[
 \frac{\eta(5/8)}{5/8}<\frac65\eta(1/10),
\]
while $\eta(t)/t$ decreases. Write $x=t^2\le25/64$.
Since $q\ge99/100$, $b\ge1/2$, and the relevant series have positive
coefficients,
\[
 \eta(t)\ge L-\frac x2-\frac{x^2}{12(1-x)},\qquad
 S(bt)\ge1+\frac{x}{12}.
\]
The entropy lower bound is positive on this interval. Thus
\begin{align*}
 K_{\rm mid}-L
 &\ge\left(L-\frac x2-\frac{x^2}{12(1-x)}\right)(1+x/12)
       +\frac{99}{200}(x+x^2/6)-L\\
 &\ge\frac{x}{1-x}
 \left(\frac{211}{4000}-\frac{381x}{4000}-\frac{43x^2}{900}\right)>0.
\end{align*}
The last polynomial decreases on this interval and equals
$152113/18432000>0$ at $x=25/64$; the substitution used $L>693/1000$.
Equation~\eqref{eq:I16-target-L} proves the product bound on this cap.

\subsection{The zero-noise cap}
Suppose $0<p\le1/100$. Appendix~\ref{app:midpoint-monotone} proves that
$K_{\rm mid}(p,t)$ decreases with $t$. It therefore suffices to use the budget
endpoint $t_\beta$, determined by $\eta(t_\beta)/t_\beta=\beta H$.
Put $\epsilon=(1-t_\beta)/2$ and $\ell=\log(1/p)\ge23/5$.
The elementary entropy remainder is
\begin{equation}\label{eq:I16-B-remainder}
 h(x)=x[\log(1/x)+B(x)],\qquad
 B(x)=1-\sum_{k\ge1}\frac{x^k}{k(k+1)},\qquad
 1-x\le B(x)\le1-x/2\le1.
\end{equation}
Set $k=5\beta/4-1/10\in[7/5,12/5]$. These bounds imply
\[
 \frac{h(kp)-\beta H}{p}
 \ge(k-\beta)(\ell+1)-k\log k-k^2p
 \ge(k-\beta)\frac{28}{5}-k\log k-\frac{k^2}{100}.
\]
The last expression is concave in $\beta$. At its endpoints it exceeds
$1561/2500$ and $44/625$, respectively, using
$\log(7/5)<17/50$ and $\log(12/5)<22/25$.
Since $h(x)/(1-2x)$ increases on $(0,1/2)$, this proves
\begin{equation}\label{eq:I16-epsilon}
 \epsilon<kp.
\end{equation}
At $t_\beta$, the exact identity
\[
 Y-2\alpha=2L+\frac{1}{2(1-p)}
 \left[\log\frac{1+(1-p)t_\beta}{1-(1-p)t_\beta}-\ell-B(p)\right]
\]
now gives a linear lower bound. Indeed, set $C_\beta=1+2k=4/5+5\beta/2$.
Then $1-(1-p)t_\beta\le pC_\beta$ and
$1+(1-p)t_\beta\ge2-pC_\beta$. The tangent bound
\[
 \log C_\beta\le\log(19/5)+\frac{25}{38}(\beta-6/5)
 <\frac{167}{125}+\frac{25}{38}(\beta-6/5)
\]
and $\log(1-pC_\beta/2)\ge-29/971$ therefore imply
\begin{equation}\label{eq:I16-Y-lower}
 Y>2\alpha+z_\beta,\qquad
 z_\beta=\frac{541}{1000}-\frac{\beta-6/5}{3}>0.
\end{equation}
For an explicit rational check, the preceding lower expression for $Y-2\alpha$ is
\[
 \frac{1386}{1000}+
 \frac{\frac{693}{1000}-\frac{167}{125}
       -\frac{25}{38}(\beta-6/5)-\frac{29}{971}-1}{2(99/100)}.
\]
Its difference from $z_\beta$ is positive at both endpoints and affine in $\beta$.
Also \eqref{eq:I16-B-remainder} gives
$\alpha\ge(\ell+1-p)/4-L>1407/2000>7/10$.

At the budget endpoint,
\[
 \frac{K_{\rm mid}(p,t_\beta)}q
 =\beta(L+\alpha)Y+\frac{Y^2}{2}.
\]
Substitute \eqref{eq:I16-Y-lower} and subtract the target divided by $q$.
The result exceeds $\alpha f_1(\beta)+f_0(\beta)$, where
\[
 f_1=2\beta\frac{693}{1000}-\frac{13}{4}+(\beta+2)z_\beta,
 \qquad
 f_0=\beta\frac{693}{1000}z_\beta+\frac{z_\beta^2}{2}-\frac{347}{500}.
\]
Both $f_1$ and $(7/10)f_1+f_0$ are concave quadratics. The first is positive
at both endpoints. The endpoint values of the second are
$33161/10^7$ and $2832997/18000000$, both greater than $1/400$.
Thus the product inequality holds with surplus $q/400$ on this entire cap.

\subsection{The compact certificate: concavity, tangents, and quartics}\label{sec:compact}
The remaining rectangle is exactly
\begin{equation}\label{eq:I16-rectangle}
 \frac1{100}\le p\le\frac9{20},\qquad 0\le t\le1.
\end{equation}
It meets both analytic caps without rounding. Each box is tested using an
entropy secant and a tangent to $S$, rather than bounding their factors
independently at opposite corners.

For $[p_l,p_h]\times[t_l,t_h]$, write $t(z)=(1-z)t_l+zt_h$, $0\le z\le1$,
and $t_0=(t_l+t_h)/2$. Let $h_-(z)$ interpolate rigorous lower bounds for
$\eta(t_l),\eta(t_h)$. Concavity of $\eta$ gives $h_-\le\eta(t(z))$.
The positive series $S(x)=\sum_{n\ge0}x^{2n}/(2n+1)$ proves convexity and
monotonicity. Enclose the value and derivative of $S((1-p_h)t)$ at $t_0$
by $[s_l,s_h]$ and $[v_l,v_h]$. The affine function with endpoint values
\[
 s_-(0)=s_l+v_h(t_l-t_0),\qquad
 s_-(1)=s_l+v_l(t_h-t_0)
\]
lies below its tangent, and hence below $S((1-p)t(z))$ throughout the box.
The derivative formula used is
\[
 \frac{\mathrm d}{\mathrm dt}S(bt)
 =\frac{(1-b^2t^2)^{-1}-S(bt)}t.
\]
All positive-leaf tests require $s_-\ge0$ at both endpoints.

For $t_l>0$, set
\[
 \beta_u=\min\left\{2,\frac{\eta(t_l)^+}{t_lH(p_l)^-}\right\};
\]
put $\beta_u=2$ when $t_l=0$. A box is entropy-infeasible if
\begin{equation}\label{eq:I16-empty-test}
 \frac{\eta(t_l)^+}{t_lH(p_l)^-}<\frac65.
\end{equation}
Superscripts indicate outward bounds, not new functions.
Otherwise every admissible $\beta$ is at most $\beta_u$.
Corollary~\ref{cor:I161-target} permits evaluation of the target at $p_h$.
One sufficient polynomial is
\begin{equation}\label{eq:I16-box-test}
 \Pi_C(z)=h_-(z)s_-(z)+\frac{q(p_l)}2t(z)^2s_-(z)^2
             -T_{\beta_u+1}(p_h)^+.
\end{equation}

A second test retains the entropy budget before bounding the product.
Since $\beta\le\eta(t)/(tH)$,
\[
 K_{\rm mid}-T_{\beta+1}
 \ge\eta(t)\left[S((1-p)t)-\frac{k(p)}t\right]
       +\frac q2t^2S((1-p)t)^2-T_1(p),
 \quad k(p)=\frac{2\alpha(p)^2}{L+\alpha(p)}.
\]
Here $k$ decreases with $p$. For $t_l>0$, let
$r(z)=(1-z)/t_l+z/t_h$ and let $k_+\ge k(p_l)$ be a rational upper bound.
Convexity of $1/t$ gives $r(z)\ge1/t(z)$. If
$d_-(z)=s_-(z)-k_+r(z)$ is nonnegative at both endpoints, another sufficient
polynomial is
\[
 \Pi_B(z)=h_-(z)d_-(z)+\frac{q(p_l)}2t(z)^2s_-(z)^2-T_1(p_h)^+.
\]
All substitutions now multiply nonnegative factors. Both polynomials have
at most degree four. Strict positivity of their five Bernstein coefficients
proves positivity on the whole box, including its boundary; the coefficient
formula is supplied in Appendix~\ref{app:geometry-cert}.

The complete certificate has 439 nodes: 219 exact midpoint bisections,
52 positive $\Pi_C$ leaves, 132 positive $\Pi_B$ leaves, and 36
entropy-infeasible leaves. All 920 positive-leaf Bernstein coefficients
are strictly positive. The smallest certified leaf margin exceeds
$7.43\times10^{-5}$; this decimal is descriptive, not a proof premise.
Each endpoint, enclosure and sign is checked with rational arithmetic and
eight logarithm-series terms. Appendix~\ref{app:replay} specifies the complete
439-character action certificate, exact coverage, and a second polynomial
reconstruction. This establishes the product lemma on the rectangle; the two
analytic caps complete its full domain.

\section{Three geometric coefficients}\label{sec:geometry}
\begin{lemma}\label{lem:I16-geometry}
The coefficients in \eqref{eq:I16-F0}--\eqref{eq:I16-G} satisfy
\begin{equation}\label{eq:I16-geometry}
 F_0>\frac1{250},\qquad F_1>\frac1{10},\qquad G>\frac1{106}
\end{equation}
throughout the parent domain in \eqref{eq:I16-domain}, independently of $R$.
The first two estimates use the six exact rational polynomial certificates
specified below and reconstructed in the replay.
\end{lemma}

\subsection{Elementary mean and gap bounds}
Put $y=(1-m)/2$. Equation~\eqref{eq:I16-B-remainder} gives
\begin{equation}\label{eq:I16-A-bounds}
 A=2\mathcal B(y),\qquad 2L\le A\le2-y\le2-V/4<2.
\end{equation}
The lower bound uses that $\mathcal B$ decreases and
$\mathcal B(1/2)=L$. Also \eqref{eq:I16-D} and the entropy series give
\begin{equation}\label{eq:I16-D-bounds}
 c_2s\le\overline D-c_1\le\frac{L-1/2}{V}.
\end{equation}
The upper bound follows by applying $\kappa(z)\le Lz^2$ to each term of
the exact binary-information identity, giving $D\le Ls/V$.
Let $x=s/V$; then
\[
 0<x\le x_*:=\min\{1/4,(1-m)/(1+m)\}.
\]
The two expressions for $x_*$ meet at $m=3/5$.

\subsection{The constant coefficient \texorpdfstring{$F_0$}{F0}}
Write $\eta=\eta(m)$ and $c_0=L-1/2$. From its definition,
\begin{align*}
 (A+x)F_0&=H_0+T(\overline D-c_1),\\
 H_0&=L(A+x)-c_0A(A+1)-x\eta/V-Ax/(2V),\\
 T&=A^2-2\eta(A+x).
\end{align*}
The sign of $T$ need not be chosen: equation~\eqref{eq:I16-D-bounds} gives
\[
 (A+x)F_0\ge\min\{H_+(x),H_-(x)\},
\]
where
\[
 H_+=H_0+T\frac{(4-3V)x}{12V^2},\qquad
 H_-=H_0+T\frac{c_0}{V}.
\]
The first is a concave quadratic in $x$; the second is affine. Subtracting
$\epsilon_0(A+x)$, with $\epsilon_0=1/250$, preserves concavity. At $x=0$,
$T=A(A-2\eta)\ge0$ by \eqref{eq:I16-A-bounds}, so their minimum is
\[
 A\{L-c_0(A+1)\}\ge A(3/2-2L)>A/10.
\]
It remains only to check both expressions at $x=x_*$. The four rational
polynomials in Appendix~\ref{app:geometry-cert} certify these endpoint
inequalities on $[0,3/5]$ and $[3/5,1)$. Concavity completes the whole
allowed interval of $x$, proving $F_0>1/250$.

\subsection{The linear coefficient \texorpdfstring{$F_1$}{F1}}
Rearrangement gives
\[
 F_1=\frac{13}{4}-\frac{2s\eta}{a+s}
 +2\left(\frac{a^2}{a+s}-V\right)\overline D-\frac{2a}{V}.
\]
The coefficient of $\overline D$ is positive: by $A\ge2L>4/3$ and
$x\le1/4$,
\[
 A^2-A-x\ge(4/3)^2-4/3-1/4=7/36>0.
\]
Use $\overline D\ge c_1+c_2s$ and multiply by $a+s$ to obtain
\begin{align*}
 (a+s)F_1\ge J_1(s)
 &=\frac94(a+s)-\frac{a^2}{V}-\frac{2as}{V}-2s\eta\\
 &\quad+2c_2s(a^2-aV-Vs).
\end{align*}
This is concave in $s$. At zero,
$J_1(0)/a=9/4-A\ge1/4>1/10$. At $s=Vx_*$, the two remaining rational
polynomials in Appendix~\ref{app:geometry-cert} certify
$J_1(s)>(a+s)/10$. Concavity of $J_1(s)-(a+s)/10$ proves $F_1>1/10$.

\subsection{The quadratic coefficient \texorpdfstring{$G$}{G}}
First omit the cap at two and put
\[
 G_{\rm raw}=\frac{V-s}{d}+1-a-V.
\]
For fixed $m$, its derivative in $d$ is $-1-V/d^2<0$, so the largest
permitted gap gives the smallest value. For $0\le m\le3/5$, use
$d=\sqrt V/2$ and the chord estimate $\sqrt{1-m^2}\ge1-5m^2/9$.
The function $a(m)$ has third derivative $-4/(1+m)<0$, hence
\[
 a\le2L+(4L-2)m+(2L-3)m^2
 <\frac75+\frac45m-\frac85m^2.
\]
Therefore
\[
 G_{\rm raw}>\frac1{10}-\frac45m+\frac{53}{30}m^2
 =\frac{53}{30}\left(m-\frac{12}{53}\right)^2+\frac1{106}.
\]
For $3/5\le m<1$, the largest gap is $1-m$, and $a\le2V$ gives
\[
 G_{\rm raw}\ge m^2+2m-a\ge3m^2+2m-2\ge\frac7{25}>\frac1{106}.
\]
If the cap in \eqref{eq:I16-beta} is active, instead use
\[
 G=3-a-V\ge3-3V+V^2/4\ge1/4>1/106.
\]
This proves the quadratic coefficient bound and completes
Lemma~\ref{lem:I16-geometry}.
% === 06_equality.tex ===
\section{Equality and the scope of the result}\label{sec:equality}
Define the profile deficit and the information deficit by
\[
\Delta^{\mathsf B}_p(f)=H(f\mid Y)-\Bprof(m),\qquad
\mathcal D_p(f)=\kappa(\rho)-I(f;Y).
\]
The master theorem and a direct subtraction give
\begin{equation}\label{eq:gap-decomposition}
\mathcal D_p(f)=\Delta^{\mathsf B}_p(f)+R\kappa(m)-\theta m^2
\ge\Delta^{\mathsf B}_p(f)+(1-L)Rm^2.
\end{equation}
Thus the stronger profile records an explicit bias penalty that the CK inequality alone does not contain. This inequality is uniform in dimension. It is not a distance-to-coordinate stability theorem by itself.

\begin{theorem}[Equality at nontrivial noise]\label{thm:equality}
For $0<p<1/2$, equality in~\eqref{eq:main-profile} holds exactly for constants and coordinate indicators or their complements. Equality in~\eqref{eq:CK} holds exactly for coordinate indicators and their complements.
\end{theorem}
\begin{proof}
The local closure is strict whenever $d>0$, by Theorem~\ref{thm:I16-blend} and the positive-weight assembly in Section~\ref{sec:blend}.

A nonconstant Boolean function that is not a signed coordinate either
has exactly two essential coordinates, in which case Lemma~\ref{lem:two-bit}
gives strictness, or has at least three. Inessential noisy coordinates are
independent of the output and may be discarded. In the latter case apply
Lemma~\ref{lem:normalized-gap} to three essential coordinates. The selected
coordinate satisfies the local quarter-variance condition. If $d>0$, strict local
closure applies.

If $d=0$, its two child means agree. They therefore supply $\overline J\ge\Bprof(m)$. Essentiality means $f_0-f_1$ is not identically zero. For $0<\rho<1$, the one-bit noise matrix has nonzero eigenvalues $1,\rho$; its finite tensor product is invertible. Hence $A-B=T_\rho(f_0-f_1)$ is nonzero for at least one observation of positive probability. Strict concavity of binary entropy then gives $\overline M>\overline J\ge\Bprof(m)$.

Constants and signed coordinates attain the profile directly. For the information inequality, constants have gap $\kappa(\rho)>0$, while signed coordinates have zero gap. All other functions have positive profile deficit, hence positive information deficit by~\eqref{eq:gap-decomposition}.
\end{proof}
At $p=0$, the profile is zero and all functions attain it; CK equality means that the output is balanced. At $p=1/2$, the profile is the output entropy and all functions attain both bounds.


\paragraph{Scope.}
The target, quantitative remainder and equality conclusions are unchanged.
No characterization of extremizing or realizable posterior laws is assumed.

\appendix
\section{Monotonicity of the midpoint product on the zero-noise cap}\label{app:midpoint-monotone}
For $0<p\le1/100$, put $b=1-p$, $c(t)=2\atanh(bt)$, and
\[
 F_p(t)=\frac{c(t)}t\eta(t)+p c(t)^2=2bK_{\rm mid}(p,t).
\]
We prove $F_p'(t)<0$ for every $0<t<1$, without a separation subdivision.
Let
\[
 S(t)=\frac{\atanh t}{t},\quad
 \mathscr A(t)=L-\frac12\ln(1-t^2),\quad G_0(t)=S(t)-\mathscr A(t).
\]
The convergent positive series and its endpoint sum give
\[
 G_0(t)=1-L-\sum_{n\ge1}\frac{t^{2n}}{2n(2n+1)}\ge0,
 \qquad \sum_{n\ge1}\frac1{2n(2n+1)}=1-L.
\]
Furthermore only the first two coefficients below are negative:
\[
 G_0(t)+\frac{S(t)}6
 =\frac76-L+\sum_{n\ge1}\frac{n-3}{6n(2n+1)}t^{2n}
 \ge\frac{187}{180}-L>\frac{61}{180}.
\]
Define $G_p=S-b^2\mathscr A-4b(1-b)$. Since $7b^2-6>0$ on this cap,
\begin{align*}
 G_p&=(7b^2-6)G_0+6(1-b^2)(G_0+S/6)-4b(1-b)\\
 &>(7b^2-6)G_0+2(1-b)^2+\frac{1-b^2}{30}>0.
\end{align*}
Direct differentiation gives the exact identity
\[
 F_p'(t)=-\frac{c(t)G_p(t)}{1-b^2t^2}
 +\frac{\eta(t)[2bt-c(t)]}{t^2(1-b^2t^2)}<0,
\]
because $c(t)\ge2bt$. All denominators are positive. This establishes the
monotonicity used only to verify the zero-noise cap of the single product lemma.

\section{The six rational geometric certificates}\label{app:geometry-cert}
This appendix specifies all rational functions whose positivity is used in
Section 8. Their numerators are reconstructed from these formulas, rather
than loaded as unexplained coefficients.

Write
\[
 \ell_- =\frac{693}{1000},\qquad
 \ell_+ =\frac{347}{500},\qquad
 \ell_-<L<\ell_+,
\]
\[
 z=\frac{1-m}{3+m},\qquad
 A_- =\frac{4(1+m)}{3+m}\left(1+\frac{z^2}{3}\right),\qquad
 A_+=A_-+\frac{4(1+m)}{3+m}\frac{z^4}{5(1-z^2)}.
\]
The identity $A=4(1+m)\atanh(z)/[(3+m)z]$ and the positive atanh series
show $A_-\le A\le A_+$. The expressions extend continuously to $m=1$.
Put $E_+=\ell_+-m^2/2\ge\eta(m)$, $c_-=\ell_--1/2$, $c_+=\ell_+-1/2$,
and $\epsilon_0=1/250$, $\epsilon_1=1/10$.

For each of the two intervals use its endpoint gap
\[
 (I,x)=([0,3/5],1/4)\quad\hbox{or}\quad
 ([3/5,1],(1-m)/(1+m)).
\]
The following lower rational functions suffice:
\begin{align}
 P_{0,+}={}&12\ell_-V^2(A_-+x)-12c_+V^2A_+(A_++1)
 -12VxE_+-6VA_+x\notag\\
 &+[A_-^2-2E_+(A_++x)](4-3V)x
 -12\epsilon_0V^2(A_++x),\label{eq:I16-polyplus}\\
 P_{0,-}={}&2V\ell_-(A_-+x)-2xE_+-A_+x
 +2c_-(1-V)A_-^2\notag\\
 &-2c_+\{VA_++2E_+(A_++x)\}
 -2\epsilon_0V(A_++x),\label{eq:I16-polyminus}\\
 P_1={}&6V\{\tfrac94(A_-+x)-A_+^2-2A_+x-2xE_+
                         -\epsilon_1(A_++x)\}\notag\\
 &+(4-3V)x(A_-^2-A_+-x).\label{eq:I16-polyone}
\end{align}
They are, respectively, lower bounds for
$12V^2[H_+-\epsilon_0(A+x)]$,
$2V[H_--\epsilon_0(A+x)]$, and the sufficient expression
$6V[J_1/V-\epsilon_1(A+x)]$ at $s=Vx$.
Every substitution respects its displayed coefficient sign. In particular,
$P_{0,-}$ separates the positive term $2c_0(1-V)A^2$ before replacing
$L$ and $A$ by upper or lower bounds.

Here is an explicit common denominator useful for independent reconstruction.
Let $T=3+m$, $Z=1-m$, and
\begin{align*}
 D_A&=30T^3,\\
 N_-&=4(1+m)(30T^2+10Z^2),\\
 N_+&=N_-+3Z^4.
\end{align*}
Then $A_-=N_-/D_A$, $A_+=N_+/D_A$. Write $x=X/Y$, with
$(X,Y)=(1,4)$ or $(1-m,1+m)$. Multiply $P_{0,+}$ and $P_1$ by
$D_A^2Y^2$, and $P_{0,-}$ by $D_A^2Y$; these multipliers are positive.
In the two indicated cases below a factor $1-m$ is then removed, which is
positive on the actual domain $m<1$.

\begin{center}
\begin{tabular}{llllr}
\toprule
Mean interval & Numerator & Degree & Removed factor & Positive coefficients\\
\midrule
$[0,3/5]$ & $P_{0,+}$ & 12 & none & 13\\
$[0,3/5]$ & $P_{0,-}$ & 9 & none & 10\\
$[0,3/5]$ & $P_1$ & 10 & none & 11\\
$[3/5,1]$ & $P_{0,+}$ & 13 & $1-m$ & 14\\
$[3/5,1]$ & $P_{0,-}$ & 10 & none & 11\\
$[3/5,1]$ & $P_1$ & 11 & $1-m$ & 12\\
\bottomrule
\end{tabular}
\end{center}
All 71 Bernstein coefficients are strictly positive, with exact rational
arithmetic. No further interval subdivision is needed for these polynomials.
For explicit reconstruction, if $P(m)=\sum_{j=0}^n p_jm^j$ on $[a,b]$, define
\[
 c_k=\sum_{j=k}^n p_j{j\choose k}a^{j-k}(b-a)^k,
 \qquad
 B_i=\sum_{k=0}^i c_k\frac{{i\choose k}}{{n\choose k}}.
\]
Then
\[
 P(a+(b-a)t)=\sum_{i=0}^n B_i{n\choose i}t^i(1-t)^{n-i}.
\]
The weights are nonnegative and sum to one on $0\le t\le1$.
The replay rebuilds every $p_j$ from \eqref{eq:I16-polyplus}--\eqref{eq:I16-polyone},
checks the exact divisions by $1-m$, computes every $B_i$, and verifies its
strict sign. The emitted JSON contains all power and Bernstein coefficients.
This proves the six endpoint inequalities used by the concavity reductions.

\section{Exact replay and provenance}\label{app:replay}
For a rational $w\in[1,2]$, put $z=(w-1)/(w+1)$. The positive logarithm
series gives the exact enclosure
\begin{equation}\label{eq:I16-log-enclosure}
 2\sum_{j=0}^{N-1}\frac{z^{2j+1}}{2j+1}
 \le\log w\le
 2\sum_{j=0}^{N-1}\frac{z^{2j+1}}{2j+1}
 +\frac{2z^{2N+1}}{(2N+1)(1-z^2)}.
\end{equation}
It follows by bounding every omitted denominator below by $2N+1$ and
summing the geometric tail. For arbitrary positive rational $x$, multiplication
or division by powers of two reduces to $[1,2)$; the contribution of $k\log2$
uses the correct endpoint order when $k<0$. For rational $x\in(0,1)$,
$h(x)=-x\log x-(1-x)\log(1-x)$ is enclosed with outward rational endpoints.
The conventions $h(0)=h(1)=0$ are exact.

The entire replay uses $N=8$. The certificate file
\texttt{product\_actions.txt} contains only the 439 preorder actions
\texttt{P}, \texttt{T}, \texttt{C}, \texttt{B}, \texttt{E}; whitespace is ignored.
The root is \eqref{eq:I16-rectangle}. The first two actions bisect the indicated
coordinate exactly at its rational midpoint, lower child first. The last three
require their respective positive-polynomial or infeasibility test. A stack
checks that no box is omitted and that there are no extra actions.
All endpoints are derived; none are stored as rounded floating-point values.
The complete action string is:
\begin{center}\begin{minipage}{0.97\textwidth}\small
\begin{verbatim}
PPPPPPTCTCPTCTCPTPBBTTPPBBPBBEETPBPBBTTPPBBPBEEETCTPCBTPPPBBPBBPPBBPBBETCTCTCPTP
BBTPPPTBBTBBPTBBTBBPPTBBTBBPTBBTBEETPBBTTPPBBPBBEETCTCPTCPTPBBTTPPBBPBBEETPPBBPB
BETPBBTPPPBBPBBPPBBPBBETCPTCPTPBBTPPPBTBBPTBBTBBPPTBBTBBTBEETPBBTTPBBEETCTPPBBPB
BTPTPPBBPBBETPPBEEEETCPTPPCCPCCTPPPBBPPBBPBBPPPBBPBBPPBBPBBETPTCPBBPBBTPPPTBBTBB
TBETPBBEETTCPCCPPTPPBBPBBTPTPBBETPBBEETPPBBTCBTPTBEEETPPTCBTCBTCPBBETTTTTCCPCCPC
TCCPPTCCTCCTPCCPCCTPPPTCCTCCTCCPTCCTCEE
\end{verbatim}
\end{minipage}\end{center}
Run \texttt{python run\_exact.py}. The command
\texttt{python product\_cert/verify\_product\_exact.py} \verb|--build| regenerates
all actions; both use the standard library only. The checker reconstructs
coverage, every leaf, the six geometric polynomials and the finite analytic
constants. A second implementation reconstructs quartics directly in Bernstein
form and checks exact agreement. Negative controls must fail. The archive
preserves ten I16 statements byte-for-byte and records the unchanged original
files' hashes. Optional high-precision tests are diagnostics, not proofs.
This remains a computer-assisted candidate without independent mathematical
approval or foundational proof-kernel verification.

\begin{thebibliography}{9}
\bibitem{CK} G. R. Kumar and T. A. Courtade.
\emph{Which Boolean Functions are Most Informative?}
arXiv:1302.2512v2 (2013).
\bibitem{I15} D. Kramer.
\emph{An entropy--variance profile for Boolean noise}, entropy-budget
iteration I15, 19 September 2026. Unpublished AI-assisted candidate.
\end{thebibliography}
\end{document}
